There's a file main/papers/src/content/bioseqdb/thesis/chapters/16_benchmarking/generated_results.tex that is generated by main/papers/scripts/generate_benchmark_results.py script.
The script extracts the results from various YAML files, takes min and creates various tables comparing performance of BioseqDB locally, on AWS machines and versus other aligners.
The yellow row is my "favourite" row marked (mpf + bounded simd + simd dp + block aligner), while first row is always minimap2 baseline.
There's also comparison (vast.ai) on machine that has better GPU than my local machine (my local machine has integrated graphisc card).

I want you now to interpret the results writing main/papers/src/content/bioseqdb/thesis/chapters/16_benchmarking/01_introduction.tex
so the vision is as follows: we first talk about results and interpretation and can sometimes mention methodology if this is relevant (for example say that "on machine with RTX XXXX we got this result"),
then user can see other stuff like methodology, machine specs, etc. the idea being that the most important part is interpretation and it should be first and it shouldn't mention all the details
like exact machine specs etc. it could be done later. What we want is just two-point list explaining shortly difference between index/no-index test cases and thats it.

The interpretation should be provided as subsequent "subsubsections" of the introduction file. Each subsubsection should discuss concrete example from some table.
I would like you to create special latex command to build small tables. The idea being that all the data is collected in many very large tables, so when we discuss 
some results we will do it one row at the time, or one column, or even some subvector of these - so few cells in row.
The results then should be always presented in the horizontal format so even if we talk about one column, it should be transposed.
The things in cell (triangle, percent comparison, A/B measurements etc.) should always stay the same format. 
Additionally the table should talk in caption what table does it refers to with \ref{}. This means that we should probably add \label to each table and each column of each table and each row of each table (not cells beacause that would be overkill) and then
this "slice" small table can refer in its caption both to source table and to the row or column of that table. We should also support a case 
to have two columns or two rows in the "slice" table (this is useful for comparison of two configurations!)

So the subsubsection looks like this:
1. We say in one sentence what is the observation without using numbers. For example "We can see that on machine M7 BioseqDB was severly slower than minimap2 for no-index case" (it's very important not to use precise numbers as I don't want to repeat myself, user can see them in the table!)
2. Here goes this "slice" table
3. We explain why the thing that we see is happening.
4. [optional] If there's more of connected examples with one root cause then you may provide more "slice" tables.

EXAMPLES (this is not the actual example but just to give you an idea what i want-  you need to provide analysis of all this cases but there might be more interesting ones - invetigation if there are others is yours job too. Additionally some of the root causes for some examples here might be the same ones and in that case this should be merged into one subsubsection):
EXAMPLE1:
All results for GPU are very slow even on vast.ai machine with great GPU because I firstly accessed CPU memory instead of GPU memory. 
That was the case for local, and on vast.ai I changed it to use GPU memory.
There are many rejections because still that makes it problematic because entire index/fragment (I don't know what is in the buffer?) doesn't likely fit in GPU.
Additionally G3SA paper describes approach where entire pipeline can be placed in one kernel so we avoid frequent host copies, but it still needs to handle cases when the buffer doesn't fit in memory.
EXAMPLE2:
The results for k=20, w=15 are slower than k=15, w=10 for baseline and other algorithms because for longer k-mers there are more possible k-mers so the has-maps are larger incresing likelihood of using SWAP?
EXAMPLE3:
blastn has usually very high startup cost compared to other tools.
EXAMPLE4:
For some reason "AWS: 2 Gbp, k=20, w=15 / index" algorithms (lets take "MPF / Scalar Scalar-DP / Block" for comparison but there are many other configs for which this holds)
work far far worse than "AWS: 4 Gbp, k=15, w=10 / index" (for this case theres several times difference!). I think this is because we have optimised kernel for 15/10 but for the other one it's some generic kernel? Need to investigate the code!!!
EXAMPLE5:
Usualy if there's little difference between A and B which means that majority of the time is likely spent on loading data/finding nachors/chaining and backtracking etc. and not on final stages of alignment MAPQ/configurations
THere's however difference observable in case of some "AWS: 2 Gbp, k=20, w=15 / index" configurations like "MPS / Scalar SIMD-DP / Block" and some others there where
A is >2X slower but B is >2X faster. This means that divergence estimation is not correctly optimised or CIGAR needs to do some copying or something that slows everything tragically (we have metnions of divergence sort problem in the minimizer chapter - go read that!)
EXAMPLE6:
This one is trivial: BioseqDB index is faster than no-index and no-index is severly slower than baseline, but this is expected because our indexing (please read main/papers/src/content/bioseqdb/thesis/chapters/12_index_storage chapter) is more sophisticated. This is acceptable as user is expected to load data once and he can wait a little bit longer for load to execute queries much faster.
EXAMPLE7:
Sometimes BioseqDB scalar alternatives for seeding and chaining are faster than SIMD (bounded SIMD + SIMD-DP vs scalar / scalar-DP) and I think this is becasue of efficient autovectorization and most of the time is spent just loading index?
EXAMPLE8:
M4 is slower in general and I think when it requests SWAP then all SWAPs are so much slower (the penatly for SWAP on M7 is lower than on M4 so anything that needs SWAPing have terrible perofrmance also you need to see specs of it and C7G beause they are solver in general but M4 is the worst!)
EXAMPLE9:
Some cases in "AWS: 4 Gbp, k=20, w=15 / index" are 2h 47 mins or above 1h - we needs to explain why is that and what happened. I think this requires investigation of the benchmarking E2E code.
Sometimes I observed that postgres was stuck on 0% CPU usage, so maybe that's the bug in the TBB pool or memory admission bug that halts everything or maybe that's the issue with SWAP?
We need deep extensive theoretical analysis of both benchmarking and BioseqDB logic to understand why this is so bad!
EXAMPLE10:
I think in cases of index we need to explain that theres no real difference between tuples like "MPF / SIMD-ref Scalar-DP / Block" and "MPF / Scalar Scalar-DP / Block" for "index" cases (for "no-index" it does matter!)
so it should be explained (I think this should be first and there should be \warn box) to explain this is because MPF plays role in how we load minimizes but seeding algorithm does not as index is laready precomputed in "index" cases.
So this should be explained as measuring the same thing twice!!!!
EXAMPLE11:
I think SIMD speedups on local are smaller than SIMD speedups on AWS (for "index" cases) because of the number of lanes? I applied the optimsiation of duplicating instructions via vector_stream in C++
but not in all cases - you need to investigate code to know for sure!
EXAMPLE12:
DNA5 on local is tragically slower (k=15,w=10) than DNA4 and that needs to do with inneficient "N" handling in kernel or something?
EXAMPLE13:
Compare few fastest tuples for AWS and local (aggregate results somehow) to tell whether it's generally better to use RMQ, SIMD or scalar for chaining and 
improtant here is to reject outliners and look carefuly only at relevant tuples/cases!
EXAMPLE14:
Do similar analysis like in EXAMPLE13 but for KSW2 and Block aligner, explain why Block or KSW2 is faster - maybe adusting block size might help for longer or smaller reads? Important to look at the code and find out! Mabe block aligner is poorly optimised or KSW2 is or maybe that some theoretical limitations like KSW2 having certain cutoff optimisation but Block not?
EXAMPLE 15:
Do similar comparison as EXAMPLE13/14 but to discuss whether MPF or MPS is faster for "index" querying and which of them is faster to build and why is that? why querying is faster what technical decisions played role here? why building is slower or faster which decisions played role here?\
EXAMPLE16:
For some reason for " AWS: 4 Gbp, k=20, w=15 / index" case "MPS / Bounded SIMD-DP / Block" is very slow on M7A for DNA5 while other cases seems to be fine on other machines.
This means that for some reason MPS there is super slow??? or this is fault of some other setting here, but I think generally MPF on M7A/M7I for DNA5 index k=20,w=15 is fast, but MPS is super super slow.
I know we do spill-sort but maybe there's no optimisation to do it more in RAM efficient way or sometrhing?
EXDAMPLE17:
minimap2 baseline for AWS: 4 Gbp, k=20, w=15 is terrible for C7G for index/no-index (indexing just reduces time by half but it's still solid 25 minutes).
Is this measured correctly? Maybe there's some problem with this instance in general?
GENERALLY:
SIMD algorithm should be the fastest, while SIMD-ref should not be because that's naive implementation of scalar. We should ivnestigate any cases where it's not the case.
Investigate  GPU slowness where G3SA declares amazing 100X speedups. Maybe GPU is better for very short sequences? We should study DOCS to find out (logs from old implemetation, paper analysis etc.)
We should ivnestigate any cases why theres sudden slowdown from one case to another and investigate any susicious behaviour like mentioned case where B is 20 seconds but A is over a miunute!
There are some very very anomalous measurements like " AWS: 2 Gbp, k=20, w=15" for "MPF / SIMD-ref SIMD-DP / Block" where suddently C7G no-index is several times faster than
index (MAKES NO SENSE THIS IS VERY LIKELY A BUG EITHER IN BENCHMARKING OR IN THE BIOSEQDB CODE OR SOME OPTIMISATION THAT GOES WRONG WAY)

Your task is to dispatch army of subagents, but before that read all chapters by yourself starting at main/papers/src/content/bioseqdb/thesis/main.tex and following all includes
(there are some legacy chapters that are not present in the final thesis). The build for some reason does not work with infinite recursion so you need to debug this as well and make sure it builds correctly.
After you are done perform at least 5 rounds of review to check if your presented argumentation is correct. Perform specialised subagent review as well to write using simple english. You ahve plenty of examples of me explaining hard concepts and it walys follow: present concept in high-level, then go into more details, discuss implications.
All things you talk about in your argumentation should refer to previous listings/figures/pseudocode/sections and subsections that I have in my work - for example when talking that this algorith case X does Y so this means Z that's why we see this results, 
we need to point user to what X is and why Y is important (by using \ref{} to point to previouis source material). It's very important to do that!!!!
Use subagents to discuss another ideas for what to describe within the measurements. You are allowed to sarch web, investigate DOCS, code, see online papers for tools I compared to determine root causes of the measurements we observe.
Do not use bioseq-build 'subagent' tool, use native subagents!!!
Make sure every "slice" table viusiaully looks similar to the original table and it's simplified i.e there arent many footnotes, it looks clean etc.
